% Generated draft: requires mathematical and source review.
\documentclass[11pt]{article}
\usepackage[margin=25mm]{geometry}
\usepackage{fontspec}
\setmainfont{Latin Modern Roman}
\usepackage{amsmath,amssymb,mathtools,unicode-math}
\setmathfont{Latin Modern Math}
\usepackage{enumitem,needspace,xurl,hyperref}
\setlength{\parindent}{0pt}
\setlength{\parskip}{5pt}
\setlength{\emergencystretch}{3em}
\newcommand{\N}{\mathbb{N}}
\newcommand{\Z}{\mathbb{Z}}
\newcommand{\Q}{\mathbb{Q}}
\newcommand{\R}{\mathbb{R}}
\newcommand{\C}{\mathbb{C}}
\newcommand{\F}{\mathbb{F}}
\newcommand{\A}{\mathbb{A}}
\newcommand{\PP}{\mathbb{P}}
\newcommand{\E}{\mathbb{E}}
\newcommand{\eps}{\varepsilon}
\newcommand{\dd}{\,\mathrm{d}}
\begin{document}
\section*{UnsolvedMath ID: 30006057}
\textbf{Source problem code:} catalogue entry
\subsection{Definitions and mathematical statement}
Let \(M\) be a connected component of the unspecified \(3\)-manifold supporting one of the explicit taut codimension-one foliations \(\mathcal F\) intended by the source. Let \(p:\widetilde M\to M\) be its universal covering map and let \(\widetilde{\mathcal F}\) be the lifted foliation. Define its leaf space by
\[
L=\widetilde M/\!\sim,\qquad x\sim y\iff x\text{ and }y\text{ belong to the same leaf of }\widetilde{\mathcal F},
\]
with the quotient topology. Deck transformations preserve the lifted foliation and therefore induce an action of \(G=\pi_1(M)\) on \(L\). The leaf space is not assumed to be Hausdorff. The foliation is called \(\mathbb R\)-covered when \(L\) is homeomorphic to \(\mathbb R\); the supplied background particularly concerns examples that are not \(\mathbb R\)-covered.

For a closed manifold, the usual closed-transversal convention calls a foliation taut if every leaf meets a closed curve transverse to the foliation. The record does not specify whether the intended manifolds are closed or have boundary, and does not supply a boundary version of tautness, regularity assumptions, or a coorientation hypothesis.

\textbf{Problem.} Understand the leaf spaces of some explicit taut foliations constructed in the literature, including the examples referred to as reference [12] in the report, sufficiently to apply the report's intended technique. The supplied edited statement describes the examples as constructed from knot diagrams and asks whether the induced fundamental-group action descends to a nontrivial orientation-preserving action on \(\mathbb R\).

If ``descends'' means a topological quotient of the leaf-space action, the corresponding precise test for a selected example is whether there exist a surjective quotient map \(q:L\to\mathbb R\) and a homomorphism
\[
\rho:G\longrightarrow\operatorname{Homeo}^{+}(\mathbb R)
\]
satisfying
\[
q(g\cdot\ell)=\rho(g)(q(\ell))\quad(g\in G,\ \ell\in L),
\qquad
\exists g\in G\ \exists t\in\mathbb R:\ \rho(g)(t)\ne t.
\]
Here \(\operatorname{Homeo}^{+}(\mathbb R)\) is the group of increasing self-homeomorphisms of \(\mathbb R\). This quotient interpretation is conditional: the supplied excerpt does not specify the required quotient map or the technique. It also does not identify the example family sufficiently to state a universal assertion, and its request concerning ``some'' examples must not be replaced by a claim that every explicit taut foliation admits such a quotient.
\subsection{Short English statement}
Analyze the lifted leaf spaces of some explicit taut foliations and determine whether their fundamental-group actions yield nontrivial actions on \(\mathbb R\) by increasing homeomorphisms. The intended examples, ambient manifolds, boundary conventions, and meaning of descent require clarification; the displayed equivariant quotient condition is a conditional formulation.
\subsection{Sources}
\begin{itemize}
\item Supplied local UnsolvedMath registry record, permanent ID \(30006057\), title \emph{Leaf Spaces of Explicit Taut Foliations}, exact source code \texttt{catalogue entry}. The record lists proposed year 2024 and supplies no named proposer.
\item Supplied report citation: \emph{MATRIX-MFO Tandem Workshop: Invariants in Low-Dimensional Topology: Combinatorics, Geometry, and Computation} (2025), Oberwolfach Reports, pp.\ 2511--2550, DOI 10.4171/owr/2024/43. Supplied source and literature link: \url{https://doi.org/10.4171/owr/2024/43}, also labeled \emph{Oberwolfach Report 43/2024}.
\item The original excerpt refers to examples in the report's reference [12], but the bibliographic entry for [12] and the surrounding explanation of the technique are not supplied.
\item Literature assessment dated exactly \texttt{2026-08-22}: supplied dated curation, not independently verified here. It labels the problem open, describes understanding explicit non-\(\mathbb R\)-covered taut-foliation leaf spaces as unresolved, states that the needed quotient action is generally difficult to obtain, and reports that no general solution was located. The supplied background expressly requests independent verification of time-sensitive status and cited claims.
\item The registry labels the edited statement \texttt{corrected\_verified} and reports that the official report was checked, the question isolated, and notation repaired. These are supplied provenance claims, not independent verification in this draft.
\end{itemize}
\end{document}
